\section{Escalas cosmol\'ogicas asociadas al ciclo de 108 estados y al operador de Perron}
\label{rm:ch:cosmo}

\subsection{Separación tipológica de las \(2\) realizaciones cosmológicas}

La construcci\'on cosmol\'ogica comprende dos magnitudes internas de
procedencia distinta. La primera se deriva del cierre peri\'odico CT108 y
determina una escala de curvatura. La segunda se deriva del cociclo de
Perron y determina una ley de diluci\'on de
memoria. Ambas usan el mismo tronco \(\APP\to\TRIT\to\TPK\), pero no tienen el
mismo dominio ni el mismo selector:

{\footnotesize
\begin{equation}
\TPK\longrightarrow\text{estado enriquecido}
\longrightarrow
\begin{cases}
\text{estructura cronol\'ogica }108\longrightarrow
R_{108}\longrightarrow(H_{108},\Lambda_{108},T_{\rm dS},S_{\rm dS}),\\
\text{cociclo de Perron}\longrightarrow
(a_n,M_n)\longrightarrow(H_\varphi,s^2\propto a^{-6}).
\end{cases}
\label{rm:eq:cosmo-two-branches}
\end{equation}
}

La igualdad eventual de unidades o de orden de magnitud no autoriza a
identificar estas realizaciones. El primer criterio de refutaci\'on del
cap\'itulo es la pretensi\'on de identificarlas sin un entrelazador que
preserve la estructura cronol\'ogica, la orientaci\'on,
memoria y carta dimensional.

\subsection{Condici\'on de periodicidad del ciclo de 108 estados}

Sea \(\ell_0=c_{\rm int}t_0\). El periodo y la frecuencia angular del retorno
son

\begin{equation}
T_{108}=108t_0,
\qquad
\omega_{108}=\frac{2\pi}{108t_0},
\qquad
R_{108}=\frac{c_{\rm int}}{\omega_{108}}
=\frac{54}{\pi}\ell_0.
\label{rm:eq:cosmo-R108}
\end{equation}

La longitud de la vuelta completa es

\begin{equation}
C_{108}=2\pi R_{108}=108\ell_0.
\label{rm:eq:cosmo-C108}
\end{equation}

Para cualquiera de las dos secciones orientadas de acci\'on
\(a\in\{\mathrm{pre},\mathrm{ret}\}\), el cuanto circular

\[
E_{108,a}=\hbar_a\omega_{108},
\qquad
m_{108,a}=E_{108,a}/c_{\rm int}^2
\]

satisface simult\'aneamente

\begin{equation}
\bar\lambda_C(m_{108,a})=R_{108},
\qquad
\lambda_C(m_{108,a})=C_{108}.
\label{rm:eq:cosmo-compton}
\end{equation}

La realizaci\'on circular, por tanto, no toma una longitud de Planck exterior.
Construye primero un radio y demuestra despu\'es que las lecturas reducida y
completa de Compton son sus dos realizaciones.

\subsection{Selector gravitatorio y relaci\'on Planck--TRIT}

El transductor gravitatorio homog\'eneo es

\begin{equation}
G_a(\theta)=\theta\frac{c_{\rm int}^5t_0^2}{\hbar_a}.
\label{rm:eq:cosmo-Gtheta}
\end{equation}

Sobre el cuanto circular, el radio gravitatorio reducido vale

\[
r_{{\rm g},a}(\theta)
=\frac{G_a(\theta)m_{108,a}}{c_{\rm int}^2}
=\theta\frac{\pi}{54}\ell_0.
\]

La condici\'on \(r_{{\rm g},a}=R_{108}\) selecciona de modo \'unico

\begin{equation}
\theta_{108}=\left(\frac{54}{\pi}\right)^2.
\label{rm:eq:cosmo-theta108}
\end{equation}

La coordenada no depende de la hoja de acci\'on: al multiplicar
\(G_a(\theta_{108})\hbar_a\), la orientaci\'on pre/retorno se cancela. De aqu\'i
resultan

\begin{align}
\ell_P&=R_{108}=\frac{54}{\pi}\ell_0,
&
t_P&=\omega_{108}^{-1}=\frac{54}{\pi}t_0,
\label{rm:eq:cosmo-planck-length-time}\\
E_P&=\hbar_a\omega_{108}
=\frac{\pi}{54}\frac{\hbar_a}{t_0},
&
E_P&=k_B\Theta_{\rm clk}\log3.
\label{rm:eq:cosmo-planck-trit}
\end{align}

La \cref{rm:eq:cosmo-planck-trit} expresa la confluencia estructural:
\(\log3\) es el contenido informacional de una decisi\'on TRIT y, en la
carta t\'ermica, convierte la temperatura cronol\'ogica en el cuanto de
Planck. La identidad no equipara el espacio de estados TRIT con el espacio
gravitatorio; demuestra que dos realizaciones funcionales determinan el
mismo escalar mediante mapas distintos.

\subsection{Realizaci\'on condicionada de la geometr\'ia de Sitter}

Supongamos ahora que el realizador Palatini--Cartan--Holst selecciona el radio
interno \(R_{108}\) como radio de un fondo de Sitter de cuatro dimensiones. La
hip\'otesis es geom\'etrica y queda visible: no se deduce de una coincidencia
num\'erica con \(H_0\). Bajo ella,

\begin{equation}
H_{108}=\frac{c_{\rm int}}{R_{108}}
=\frac{\pi}{54t_0},
\qquad
\Lambda_{108}=\frac{3}{R_{108}^2}
=\frac{\pi^2}{972\ell_0^2}.
\label{rm:eq:cosmo-ds-H-Lambda}
\end{equation}

La temperatura de Gibbons--Hawking y la entrop\'ia de horizonte quedan

\begin{align}
k_BT_{\rm dS}
&=\frac{\hbar_a H_{108}}{2\pi}
=\frac{\hbar_a}{108t_0},
&
\frac{T_{\rm dS}}{\Theta_{\rm clk}}
&=\frac{\log3}{2\pi},
\label{rm:eq:cosmo-ds-temperature}\\
\frac{S_{\rm dS}}{k_B}
&=\frac{\pi R_{108}^2}{\ell_P^2}=\pi.
\label{rm:eq:cosmo-ds-entropy}
\end{align}

La \(\pi\) de \cref{rm:eq:cosmo-ds-entropy} no se introduce como valor de una
entrop\'ia a reproducir: aparece porque el selector circular hace
\(R_{108}=\ell_P\), y el \'area de la esfera de horizonte contiene el factor
geom\'etrico \(4\pi\).

\begin{proposition}[Familia adimensional CT108--de Sitter]
Bajo la realizaci\'on de Sitter declarada, el cu\'adruple
\[
\left(
H_{108}t_0,
\Lambda_{108}\ell_0^2,
\frac{T_{\rm dS}}{\Theta_{\rm clk}},
\frac{S_{\rm dS}}{k_B}
\right)
=\left(
\frac\pi{54},
\frac{\pi^2}{972},
\frac{\log3}{2\pi},
\pi
\right)
\]
es una salida adimensional exacta.
\end{proposition}

\begin{falsifier}
La realizaci\'on de Sitter queda invalidada si el fondo seleccionado no
satisface \(R_{\mu\nu}=\Lambda_{108}g_{\mu\nu}\), si el realizador PCH determina
un radio distinto de \(R_{108}\), o si las cartas usadas en
\cref{rm:eq:cosmo-ds-H-Lambda,rm:eq:cosmo-ds-temperature} no comparten
\((c_{\rm int},t_0,\hbar_a)\). Ninguno de esos fallos altera el teorema
circular \cref{rm:eq:cosmo-R108,rm:eq:cosmo-theta108}.
\end{falsifier}
